% 原书第117页。按原表文字重排；允许跨页，重复表头。
\begingroup\small\setlength{\tabcolsep}{4pt}\renewcommand{\arraystretch}{1.18}
\setlength{\LTpost}{3pt}\setlength{\parskip}{0pt}
\begin{longtable}{|p{3em}|p{\dimexpr.5\linewidth-1.5em-3\tabcolsep-2\arrayrulewidth\relax}|p{\dimexpr.5\linewidth-1.5em-3\tabcolsep-2\arrayrulewidth\relax}|}
\caption{三位皇子联名奏折与《实录》关于太子废立前后事记载对比}\label{tab:comparison}\\
\hline
 & 皇三子等雍正四年奏折 & 《圣祖仁皇帝实录》\\\hline\endfirsthead
\hline & 皇三子等雍正四年奏折 & 《圣祖仁皇帝实录》\\\hline\endhead
\endfoot
大皇子 & 大阿哥奏云：二阿哥行止卑污，大失人心。山东相面人张明德曾相八阿哥，后必大贵。 & 允禔奏：允礽所行卑污、大失人心。相面人张明德曾相允禩，后必大贵。\\\hline
 & 今果欲诛二阿哥，不必出自皇父之手。朕闻之殊觉惊异，因密止大阿哥之谋。 & 今欲诛允礽，不必出自皇父之手。言至此，朕为之惊异。\\\hline
 & 伊本系一暴悍愚昧、不明大义之人，若果同八阿哥聚集贼徒，谋害二阿哥，彼时岂暇计及于朕躬有碍否耶？ & 朕思允禔为人凶顽愚昧，不知义理。倘果同允禩聚集党羽，杀害允礽，其时但知逞其凶恶，岂暇计及于朕躬有碍否耶？\\\hline
 & 此皆不明大义、不守君臣定分之人，实国家之乱臣贼子，天理国法断皆不容宽宥者也。 & 似此不谙君臣大义、不念父子至情之人，洵为乱臣贼子，天理国法皆所不容者也。\\\hline
皇八子 & 八阿哥自幼赋性阴险，存心僭越，其纠合党类，妄行作乱者有。 & 允禩自幼性奸心妄。其纠合党类、妄行作乱者有。\\\hline
 & 允禩亦甚受制于其妻。 & 再允禩素受制于妻。\\\hline
 & 其将二阿哥之事如何措置之处，悉由朕独断，谁敢多言？春秋大义，臣子无将，将则必诛。 & 废皇太子后，允禔曾奏称允禩好。春秋之义，人臣无将，将则必诛。\\\hline
 & 今众阿哥等惟当念朕为君父，朕如何降旨，尔等即如何遵行，方为臣子之所当然。尔等若不如此存心，日后朕躬考终，必至将朕躬置于乾清宫，而尔等束甲相争耳。 & 众阿哥当思朕为君父，朕如何降旨，尔等即如何遵行，始是为臣子之正理。尔等若不如此存心，日后朕躬考终，必至将朕躬置乾清宫内，尔等束甲相争耳。\\\hline
 & 允禩乃辛者库贱妇所出。 & 其母家亦甚微贱。\\\hline
\end{longtable}
\noindent（以上联名奏折内容出自：《世宗宪皇帝上谕八旗》，《四库全书》本，卷4，雍正四年二月初五日。）
\par\endgroup
